\begin{proof}[Proof of \cref{obj:thm:sharp-local-linear-constant-and-representers}]
\begin{enumerate}
\item Fix \(t\) and write \(n=n_t\), \(d=d_t\), and \(m=m_t\).  On the complete-block graph, every active block is a full neighborhood of size \(d\), and \(n=md\).  The canonical block-local weights
\[
w^0_{i,t}(Z)=g_d(Z_{b(i)})
\]
belong to \(\mathcal E_t^{\mathrm{loc,lin}}\).  The block program identities in \cref{obj:lem:block-energy-representer}, applied block by block, give value \(A_d\) per unit score and therefore \(d^2A_d\) per block in the block-extremal criterion.  Thus every admissible local-linear rule has worst-case risk at least
\[
\frac{B^2}{n^2}\sum_{b=1}^{m}d^2A_d
 = \frac{B^2A_d}{m},
\]
while the canonical rule attains this value.  Since \(n=md\),
\[
\frac{B^2A_d}{m}=\frac{B^2dA_d}{n}.
\]
This proves the finite minimax identity. % lean: locLinMinimaxRisk_exact

\item For a fixed \(w\in\mathcal E_t^{\mathrm{loc,lin}}\), the squared error decomposes across complete blocks.  Maximizing the coefficient signs and low-order monomials independently on each block gives
\[
\sup_{c\in\mathcal M_t(G_t)}
\mathbb E_Z[(\widehat\tau_{w,t}-\tau_{n_t})^2]
=
\frac{B^2}{n_t^2}\sum_{b=1}^{m_t}\Psi_{b,t}(w),
\]
where the finite maximum in \(\Psi_{b,t}\) is the same block-extreme functional after translating the sign coding to \(\varepsilon_i\in\{-1,1\}\) and writing block monomials as \(Z_{S_i}\). % lean: locLinWorstRisk_exact_blockExtremal

\item The per-block lower bound follows from the all-empty monomial candidate and the block score identity.  Let \(g\) be the common complete-block score on \(V_b\).  For every \(i\in V_b\), blockwise unbiasedness gives
\[
\mathbb E_Z[w_{i,t}(Z_b)g(Z_b)]=A_{d_t},
\qquad
\mathbb E_Z[g(Z_b)^2]=A_{d_t}.
\]
Hence
\[
0\le
\mathbb E_Z\!\left[
\left\{\sum_{i\in V_b}w_{i,t}(Z_b)-d_t g(Z_b)\right\}^2
\right]
=
\mathbb E_Z\!\left[
\left\{\sum_{i\in V_b}w_{i,t}(Z_b)\right\}^2
\right]
-d_t^2A_{d_t}.
\]
The all-empty choice in \(\Psi_{b,t}\) contains the first expectation on the right, so
\[
d_t^2A_{d_t}\le \Psi_{b,t}(w).
\]
% lean: completeBlock_blockExtremal_lower

\item Define
\[
Y_t:=
\frac{1}{m_t d_t^2 A_{d_t}}
\sum_{b=1}^{m_t}\{\Psi_{b,t}(w_t)-d_t^2A_{d_t}\}.
\]
The preceding finite identity gives, for every \(t\),
\[
\frac{
\sup_{c\in\mathcal M_t(G_t)}
\mathbb E_Z[(\widehat\tau_{w_t,t}-\tau_{n_t})^2]
}{
B^2A_{d_t}/m_t
}
=1+Y_t.
\]
Therefore convergence of the risk ratio to \(1\) is equivalent to convergence of \(Y_t\) to \(0\), by subtracting and adding the constant sequence \(1\). % lean: riskRatio_tendsto_iff_normalizedBlockExcess

\item Let
\[
X_t:=
\frac{1}{n_tA_{d_t}}
\sum_{i\in V_{n_t}}
\mathbb E_Z\!\left[
\{w_{i,t}(Z_{b(i)})-g_{d_t}(Z_{b(i)})\}^2
\right].
\]
The block-excess sum is nonnegative by the lower bound above.  The blockwise distance inequality gives
\[
Y_t\le 2\sqrt{X_t}+X_t,
\]
after using \(n_t=m_td_t\) and simplifying
\[
\frac{
2\sqrt{d_t^2A_{d_t}}\sqrt{m_t}\sqrt{d_tE_t}
+d_tE_t
}{
m_td_t^2A_{d_t}
}
=
2\sqrt{\frac{E_t}{n_tA_{d_t}}}
+\frac{E_t}{n_tA_{d_t}}.
\]
Thus \(X_t\to0\) implies \(Y_t\to0\) by squeezing, and the identity in the previous step gives the asserted convergence of the risk ratio. % lean: normalizedWeightDistance_tendsto_implies_riskRatio

\item Assume \(d_t\to\infty\).  For all sufficiently large \(t\), \(\beta<d_t\).  Define \(w_t\) by adding to the canonical weights a centered full-block perturbation on one distinguished unit in each block:
\[
w_{i,t}(Z)
=
g_{d_t}(Z_{b(i)})
+
\mathbf 1\{i=i_b^\star\}
\left\{
\frac{2d_tA_{d_t}}{(p(1-p))^{d_t}}
\right\}^{1/2}
\prod_{j\in V_{b(i)}}(Z_j-p).
\]
The full-block centered monomial has order \(d_t\), so it is orthogonal under \cref{obj:ass:bernoulli-design} to every monomial of order at most \(\beta\); consequently these weights remain in \(\mathcal E_t^{\mathrm{loc,lin}}\).  The exact identity above gives the lower bound \(1\) for the risk ratio.  The witness block bound gives, eventually,
\[
\Psi_{b,t}(w_t)
\le
d_t^2A_{d_t}
+2\sqrt{d_t^2A_{d_t}}\sqrt{2d_tA_{d_t}}
+2d_tA_{d_t},
\]
and therefore
\[
\frac{
\sup_{c\in\mathcal M_t(G_t)}
\mathbb E_Z[(\widehat\tau_{w_t,t}-\tau_{n_t})^2]
}{
B^2A_{d_t}/m_t
}
\le
1+2\sqrt{\frac{2}{d_t}}+\frac{2}{d_t}.
\]
Since \(d_t\to\infty\), squeezing proves the displayed risk-ratio convergence for this sequence. % lean: orthogonalWitness_riskRatio_tendsto_one

\item For the same sufficiently large \(t\), the perturbation has squared energy \(2d_tA_{d_t}\) on the distinguished unit of each block and zero perturbation energy on the other units.  Hence
\[
\sum_{i\in V_{n_t}}
\mathbb E_Z\!\left[
\{w_{i,t}(Z_{b(i)})-g_{d_t}(Z_{b(i)})\}^2
\right]
=
m_t(2d_tA_{d_t}),
\]
and division by \(n_tA_{d_t}=m_td_tA_{d_t}\) gives the value \(2\).  If \(\pi\) preserves block membership, the complete-block score is invariant under the induced within-block relabeling, and the full-block centered perturbation has the same squared energy after the relabeling.  The same calculation gives the relabeled normalized distance \(2\). % lean: orthogonalWitness_normalized_relabel_distance_eq_two
\end{enumerate}
\end{proof}

\begin{proof}[Proof of \cref{obj:thm:fair-coin-energy-frontier}]
\begin{enumerate}
\item Fix admissible \(n,d,\beta,B\).  For every \(r\ge 1\),
\[
\Delta_r(1/2)
=
(1/2)^r-(-1/2)^r
=
2^{-r}\{1-(-1)^r\}.
\]
Thus \(\Delta_r(1/2)=2^{1-r}\) when \(r\) is odd, and \(\Delta_r(1/2)=0\) when \(r\) is even. % lean: fairCoinContrast

\item Expanding \(A_d\) from \cref{obj:def:block-score-energy} at \(p=1/2\) gives
\[
A_d
=
\sum_{r=1}^{\bar\beta_d}
\binom{d}{r}
\frac{\Delta_r(1/2)^2}{((1/2)(1-1/2))^r}.
\]
The preceding calculation makes the even terms vanish.  For odd \(r\),
\[
\frac{\Delta_r(1/2)^2}{((1/2)(1-1/2))^r}
=
\frac{(2^{1-r})^2}{(1/4)^r}
=
4.
\]
Therefore
\[
A_d
=
4\sum_{\substack{1\le r\le \bar\beta_d\\ r\text{ odd}}}
\binom{d}{r}.
\] % lean: fairCoinEnergy
The exposed set in \cref{obj:def:exposed-order} is exactly the set of odd integers in \(\{1,\ldots,\bar\beta_d\}\).  Hence \(k_\star(d,\beta,1/2)\) is its maximum, with the stated zero value when the set is empty. % lean: fairCoinKStar

\item Apply \cref{obj:thm:bounded-outcome-degree-frontier} with \(\beta=1\) and \(p=1/2\).  It supplies constants \(a,C\) with \(0<a\le C\).  Set
\[
c_{\mathrm{lower}}:=a,
\qquad
c_{\mathrm{upper}}:=4C .
\]
Since \(C>0\), these constants satisfy \(0<c_{\mathrm{lower}}\le c_{\mathrm{upper}}\).  For fixed admissible \(n,d,B\), use the common fair Bernoulli product design on \(V_n\), which satisfies \cref{obj:ass:bernoulli-design}. % lean: bounded_outcome_degree_frontier

\item The first-order interaction energy follows from the fair-coin energy identity.  Since \(d\ge1\), \(\bar\beta_d=\min\{1,d\}=1\), and the odd orders in \(\{1\}\) consist only of \(1\).  Hence
\[
A_d
=
4\binom d1
=
4d.
\]
Put \(x=d^2/n\).  Then \(x\ge0\), and
\[
\frac{dA_d}{n}
=
\frac{d(4d)}{n}
=
4x.
\]
Moreover \(\min\{1,x\}\le \min\{1,4x\}\). % lean: henergy, henergyArg, hminLower

\item The lower bound from \cref{obj:thm:bounded-outcome-degree-frontier} gives
\[
aB^2\min\!\left\{1,\frac{dA_d}{n}\right\}
\le
R_{n,\ell_1}^{\star}(d,1,1/2,B),
\]
and the same result identifies
\[
R_{n,\ell_1}^{\star}(d,1,1/2,B)
=
R_n^\star(d,1,1/2,B).
\]
Using \(\min\{1,x\}\le\min\{1,4x\}\) and \(dA_d/n=4x\) yields
\[
c_{\mathrm{lower}}B^2\min\!\left\{1,\frac{d^2}{n}\right\}
\le
R_n^\star(d,1,1/2,B).
\] % lean: hminLower, henergyArg, hlower, hriskEq, hlower'
The class comparison in \cref{obj:thm:bounded-outcome-degree-frontier} supplies
\[
R_{n,\ell_1}^{\star}(d,1,1/2,B)
\le
R_{n,\infty}^{\star}(d,1,1/2,B).
\] % lean: hclassMono
For the upper bound, the same result bounds the bounded-outcome minimax risk by the clipped bounded-outcome SNIPE worst risk and then by
\[
CB^2\min\!\left\{1,\frac{dA_d}{n}\right\}
=
CB^2\min\{1,4x\}.
\]
If \(x\le1\), then either \(4x\le1\), in which case \(\min\{1,4x\}=4x\), or \(4x>1\), in which case \(\min\{1,4x\}=1\le4x\).  If \(x>1\), then \(\min\{1,4x\}\le1\le4\min\{1,x\}\).  Thus
\[
\min\{1,4x\}\le4\min\{1,x\},
\]
and consequently
\[
R_{n,\infty}^{\star}(d,1,1/2,B)
\le
c_{\mathrm{upper}}B^2\min\!\left\{1,\frac{d^2}{n}\right\}.
\] % lean: hbddToWorst, hworstUpper, henergyArg, hupper'
Combining the displayed identities and inequalities gives the full minimax chain.

\item Finally assume \(d\mid n\).  The exact complete-block identities from \cref{obj:thm:bounded-outcome-degree-frontier} give, for the unprojected SNIPE estimator,
\[
\sup_{\mathcal M_{n,d,1}(B)}
\mathbb E\!\left[\bigl(\widehat\tau_n^{\mathrm{SNIPE}}-\tau_n\bigr)^2\right]
=
\frac{B^2A_d}{n/d}
=
\frac{B^2dA_d}{n},
\]
and the same identity over \(\mathcal M_{n,d,1}^{\infty}(B)\).  Substituting \(A_d=4d\) gives both asserted values,
\[
\frac{B^2dA_d}{n}
=
\frac{4B^2d^2}{n}.
\]
% lean: fair_coin_energy_frontier
\end{enumerate}
\end{proof}

\begin{proof}[Proof of \cref{obj:thm:degree-frontier}]
\begin{enumerate}
\item Apply \cref{obj:thm:bounded-outcome-degree-frontier} with the fixed \(\beta\) and \(p\).  It gives constants \(a,b\) with \(0<a\le b\) and, for every admissible \(n,d,B\), the block-energy lower bound, the projected-SNIPE upper bound, the unprojected-SNIPE upper bound, the local-energy comparison, the exposed-order comparison
\[
a\,d\binom{d}{k_\star(d,\beta,p)}
\le dA_d
\le b\,d\binom{d}{k_\star(d,\beta,p)},
\]
and the block-mixture conclusions.  Set
\[
c_{\beta,p}:=\min\{1,a^2\},
\qquad
C_{\beta,p}:=\max\{1,b^2\}.
\]
Then \(0<c_{\beta,p}\le C_{\beta,p}\). % lean: degree_frontier

\item Fix \(n,d\in\mathbb N\) and \(B\in\mathbb R\) with \(n\ge1\), \(d\le n\), and \(B\ge0\), and write
\[
x:=\frac{d\binom{d}{k_\star(d,\beta,p)}}{n},
\qquad
y:=\frac{dA_d}{n}.
\]
The common-\(p\) product Bernoulli design satisfies \cref{obj:ass:bernoulli-design}.  Since \(n\ge1\), the exposed-order comparison from the preceding item gives
\[
a x\le y\le b x,
\]
and hence
\[
\min\{1,ax\}\le \min\{1,y\}\le \min\{1,bx\}.
\]
Also \(\min\{1,a^2\}\le a\): if \(a\le1\), then \(\min\{1,a^2\}\le a^2\le a\), while if \(a>1\), then \(\min\{1,a^2\}\le1\le a\).  Splitting on \(x\le1\) and \(ax\le1\), and using this comparison in the saturated branches, gives
\[
\min\{1,a^2\}\min\{1,x\}
\le
a\min\{1,ax\}.
\]
Similarly \(b\le\max\{1,b^2\}\): if \(b\le1\), then \(b\le1\le\max\{1,b^2\}\), while if \(b>1\), then \(b\le b^2\le\max\{1,b^2\}\).  Splitting on \(x\le1\) and \(bx\le1\) gives
\[
b\min\{1,bx\}
\le
\max\{1,b^2\}\min\{1,x\}.
\]
In the branch \(x>1\) and \(bx\le1\), the needed bound is obtained by first deriving \(b<1\): otherwise \(b\ge1\) would imply \(x\le bx\), contradicting \(x>1\) and \(bx\le1\). % lean: degree_frontier

\item The lower bound follows by chaining the previous inequalities with the block-energy lower bound supplied by \cref{obj:thm:bounded-outcome-degree-frontier} and the identification of the coefficient-mass minimax risk in \cref{obj:def:minimax-risk}:
\[
\begin{aligned}
c_{\beta,p}B^2\min\{1,x\}
&=B^2\min\{1,a^2\}\min\{1,x\} \\
&\le B^2a\min\{1,ax\} \\
&\le B^2a\min\{1,y\} \\
&\le R_n^\star(d,\beta,p,B).
\end{aligned}
\] % lean: degree_frontier

\item The projected-SNIPE upper bound at the block-energy scale from \cref{obj:thm:bounded-outcome-degree-frontier}, together with the upper rescaling above, gives
\[
\begin{aligned}
\sup_{(G,c)\in\mathcal M_{n,d,\beta}(B)}
\mathbb E\!\left[
  \bigl(\widehat\tau_n^{\mathrm{up}}-\tau_n\bigr)^2
\right]
&\le bB^2\min\{1,y\} \\
&\le bB^2\min\{1,bx\} \\
&\le C_{\beta,p}B^2\min\{1,x\}.
\end{aligned}
\] % lean: degree_frontier

\item The estimator \(\widehat\tau_n^{\mathrm{up}}\) is an admissible competitor for the minimax problem.  The finite-population estimator is measurable, and \cref{obj:ass:bounded-coefficient-mass} gives, for every model in \(\mathcal M_{n,d,\beta}(B)\),
\[
|Y_i(z)|\le
\sum_{S\subseteq N_i}|c_{i,S}|\le B.
\]
Therefore
\[
|\tau_n|
\le
\frac1n\sum_{i\in V_n}\bigl(|Y_i(\mathbf 1)|+|Y_i(\mathbf 0)|\bigr)
\le 2B.
\]
The projection defining \(\widehat\tau_n^{\mathrm{up}}\) places the estimator in \([-B,B]\), so its pointwise squared loss is bounded by \(9B^2\), and hence its risk is uniformly bounded over the class.  Since \(n\ge1\), the empty graph with the zero coefficient schedule supplies a member of \(\mathcal M_{n,d,\beta}(B)\).  The definition of \(R_n^\star(d,\beta,p,B)\) as the infimum over admissible estimators gives
\[
R_n^\star(d,\beta,p,B)
\le
\sup_{(G,c)\in\mathcal M_{n,d,\beta}(B)}
\mathbb E\!\left[
  \bigl(\widehat\tau_n^{\mathrm{up}}-\tau_n\bigr)^2
\right].
\] % lean: snipeClipped_admissible, minimaxRisk_le_worstRisk_of_admissible

\item The unprojected SNIPE bound is obtained from the block-energy bound in \cref{obj:thm:bounded-outcome-degree-frontier}:
\[
\begin{aligned}
\sup_{(G,c)\in\mathcal M_{n,d,\beta}(B)}
\mathbb E\!\left[
  \bigl(\widehat\tau_n^{\mathrm{SNIPE}}-\tau_n\bigr)^2
\right]
&\le B^2\,\frac{dA_d}{n} \\
&\le B^2\,b\,x \\
&\le C_{\beta,p}B^2\,x.
\end{aligned}
\] % lean: degree_frontier

\item The score-energy assertion is exactly the local-energy comparison supplied by \cref{obj:thm:bounded-outcome-degree-frontier}: for every \((G,c)\in\mathcal M_{n,d,\beta}(B)\) and every \(i\in V_n\),
\[
A_i\le A_d.
\] % lean: degree_frontier

\item For every \((G,c)\in\mathcal M_{n,d,\beta}(B)\), every \(i\in V_n\), and every order \(r\) with \(1\le r\), the overlap-count assertion is \cref{obj:lem:overlap-count} applied to the graph \(G\), the degree bound \(d\), the order \(r\), and the unit \(i\):
\[
\sum_{\ell\in V_n}\binom{|N_i\cap N_\ell|}{r}
=
\sum_{\substack{S\subseteq N_i\\ |S|=r}}
\#\{\ell\in V_n:S\subseteq N_\ell\}
\le
d\binom{|N_i|}{r}
\le
d\binom{d}{r}.
\] % lean: overlap_count_le

\item Assume \(B>0\) and \(d\ge1\).  The block-mixture component of \cref{obj:thm:bounded-outcome-degree-frontier}, with \(m=\lfloor n/d\rfloor\), \(\rho\) the active population fraction, and \(\delta\) the least-favourable perturbation amplitude, gives the two complete-block models \(M_+\) and \(M_-\) for every \(U=(U_1,\ldots,U_m)\) with \(|U_b|\le B/2\).  Their graph is the complete-block graph, their coefficient schedules are the corresponding \(+1\) and \(-1\) block schedules, and
\[
\tau_n(M_+)=\rho\delta,
\qquad
\tau_n(M_-)=-\rho\delta.
\]
The same component gives
\[
H^2(\Pi_+,\Pi_-)
\le
\frac{4\pi^2m\delta^2}{B^2A_d},
\qquad
\frac12\le\rho\le1,
\qquad
\frac{A_d}{m}=\frac{dA_d/n}{\rho}.
\]
This proves the block-construction assertion and completes the proof. % lean: degree_frontier
\end{enumerate}
\end{proof}

\begin{proof}[Proof of \cref{obj:thm:bounded-outcome-degree-frontier}]
\begin{enumerate}
\item By \cref{obj:lem:block-energy-representer}, choose constants
\(c_1,c_2,H\) depending only on \((\beta,p)\), with
\(0<c_1\le c_2\), and with
\[
c_1\binom{d}{k_\star(d,\beta,p)}
\le A_d
\le
c_2\binom{d}{k_\star(d,\beta,p)}
\qquad(d\ge 1).
\]
Let \(H_{\beta,p}\) denote the corresponding uniform raw-representer mass
supremum and set
\[
\kappa
=
\min\!\left\{(2H_{\beta,p})^{-1},(4\pi)^{-1}\right\},
\qquad
c_{\beta,p}
=
\min\!\left\{c_1,\frac{\kappa^2}{16}\right\},
\qquad
C_{\beta,p}
=
\max\{16,c_2\}.
\]
The representer-mass positivity and \(\pi>0\) give \(\kappa>0\), hence
\(c_{\beta,p}>0\), and
\[
c_{\beta,p}\le c_1\le c_2\le C_{\beta,p}.
\]
% lean: blockEnergy_representer

\item Fix \(n,D,d,B\) satisfying the hypotheses. The degree-index condition gives
\(d\le n\). The coefficient-mass envelope bounds every potential outcome by
\(B\) on the same carrier, so
\(\mathcal M_{n,d,\beta}(B)\subseteq
\mathcal M_{n,d,\beta}^{\infty}(B)\) carrier by carrier. If \(B>0\) and
\(d\ge 1\), the one-loop schedule with coefficients \(B\) and \(-2B\) has
uniformly bounded outcomes and coefficient mass \(3B>B\), giving strict
containment. The equality
\(R_{n,\ell_1}^{\star}(d,\beta,p,B)=R_n^{\star}(d,\beta,p,B)\) is the
definition of the coefficient-class minimax risk used here, and the carrier
inclusion gives
\[
R_{n,\ell_1}^{\star}(d,\beta,p,B)
\le R_{n,\infty}^{\star}(d,\beta,p,B).
\]
% lean: modelClassIncluded_of_coeffMass

\item The scaled block-prior lower bound gives
\[
c_{\beta,p}B^2\min\!\left\{1,\frac{dA_d}{n}\right\}
\le
R_{n,\ell_1}^{\star}(d,\beta,p,B).
\]
For \(d\ge1\) and \(B>0\), this is obtained with
\(m=\lfloor n/d\rfloor\), \(\rho=md/n\),
\(t=\min\{1,\sqrt{A_d/m}\}\), and
\(\delta=B\kappa t\). The choice
\(\kappa\le(4\pi)^{-1}\) makes the squared Hellinger distance of the two block
priors at most \(1/4\). Le Cam's block-prior lower bound then gives
\((\rho\delta)^2/4\), and the identities
\[
\frac{dA_d}{n}=\rho\frac{A_d}{m},
\qquad
\frac12\le\rho\le1,
\qquad
t^2=\min\!\left\{1,\frac{A_d}{m}\right\}
\]
convert this to the displayed frontier lower bound. For \(B=0\), and also for
\(d=0\), the left side is \(0\), and the minimax risk is nonnegative.
% lean: minimaxRisk_scaled_block_lower

\item The unclipped SNIPE variance calculation, with the local dependence
aggregation controlled by \cref{obj:lem:overlap-count}, yields on both model
classes
\[
\sup \mathbb E_Z\!\left[
(\widehat\tau_n^{\mathrm{SNIPE}}-\tau_n)^2
\right]
\le
B^2\frac{dA_d}{n}.
\]
For each model in the coefficient-mass class, and for each model in the
uniformly bounded-outcome class, the product Bernoulli design assumption
identifies \(D\) with the common-\(p\) Bernoulli law, and the SNIPE moment
identity gives unbiasedness for \(\tau_n\).
% lean: worstRisk_snipe_le

\item Projection to \([-B,B]\) on the coefficient-mass class and to
\([-2B,2B]\) on the uniformly bounded-outcome class can only improve the
modelwise squared error relative to the corresponding target envelopes, while
the trivial saturated bounds are \(4B^2\) and \(16B^2\). Combining these with
the unclipped risk bounds gives
\[
\sup_{\mathcal M_{n,d,\beta}(B)}
\mathbb E_Z\!\left[
(\widehat\tau_n^{\mathrm{up}}-\tau_n)^2
\right]
\le
4B^2\min\!\left\{1,\frac{dA_d}{n}\right\},
\]
and
\[
\sup_{\mathcal M_{n,d,\beta}^{\infty}(B)}
\mathbb E_Z\!\left[
(\widehat\tau_{n,\infty}^{\mathrm{up}}-\tau_n)^2
\right]
\le
16B^2\min\!\left\{1,\frac{dA_d}{n}\right\}.
\]
Since \(4\le16\le C_{\beta,p}\), the coefficient-class clipped bound and the
bounded-outcome clipped bound have the theorem's constant \(C_{\beta,p}\). The
minimax risk over the bounded-outcome class is bounded above by the risk of the
bounded-outcome clipped SNIPE estimator, giving the displayed minimax upper
chain.
% lean: worstRiskBdd_clipped_le_min

\item For every unit \(i\), the exact local score energy is bounded by the
complete-block energy:
\[
A_i\le A_d.
\]
The same estimate applies on both model classes because it depends only on the
degree bound of the carrier graph. Multiplying the energy scale from
\cref{obj:lem:block-energy-representer} by \(d\) and using
\(c_{\beta,p}\le c_1\) and \(c_2\le C_{\beta,p}\) gives, for \(d\ge1\),
\[
c_{\beta,p}d\binom{d}{k_\star(d,\beta,p)}
\le dA_d
\le
C_{\beta,p}d\binom{d}{k_\star(d,\beta,p)}.
\]
For \(d=0\), both sides reduce by the zero-degree convention.
% lean: localEnergy_le_blockEnergy

\item Assume \(d\ge1\) and \(d\mid n\), and put \(m=n/d\). On the complete
directed \(d\)-block graph with loops, first consider \(B>0\). The canonical
local-linear weights agree with the SNIPE score on each block. The block
extremal criterion equals \(d^2A_d\) on every block, so summing over \(m\)
blocks and normalizing yields
\[
\frac{B^2A_d}{m}.
\]
The fixed-graph coefficient-class risk equals this value. The global
coefficient-class risk is bracketed between that fixed-graph risk and the
unclipped SNIPE upper bound at the same value; the fixed-graph
bounded-outcome risk and the global bounded-outcome risk are obtained by the
same comparison, using the carrier-preserving inclusion of the coefficient
class in the bounded-outcome class. If \(B=0\), the target
\(B^2A_d/m\) is zero; the zero schedule gives a nonnegative risk in the
fixed-graph coefficient class, and the SNIPE upper bounds and fixed-to-global
comparisons force the fixed-graph and global risks in both model classes to be
zero. Thus, for all \(B\ge0\), the fixed-graph coefficient-class risk, the
fixed-graph bounded-outcome risk, and the corresponding global worst risks all
equal \(B^2A_d/m\). Since \(n=md\),
\[
\frac{B^2A_d}{m}=B^2\frac{dA_d}{n}.
\]
% lean: completeBlock_all_risks_exact

\item Finally suppose \(B>0\) and \(d\ge1\), set
\(m=\lfloor n/d\rfloor\), \(\rho=md/n\), and let
\(\delta\) be the perturbation amplitude. For every
\(U=(U_1,\ldots,U_m)\) with \(|U_b|\le B/2\), the two signed block schedules
are members of \(\mathcal M_{n,d,\beta}(B)\) on the complete-block graph. Their
targets are
\[
\rho\delta
\qquad\text{and}\qquad
-\rho\delta .
\]
The block-prior Hellinger calculation gives
\[
H^2(\Pi_{+},\Pi_{-})
\le
\frac{4\pi^2m\delta^2}{B^2A_d}.
\]
The active-share arithmetic gives
\[
\frac12\le\rho\le1,
\qquad
\frac{A_d}{m}
=
\frac{dA_d/n}{\rho}.
\]
These are exactly the least-favourable family conclusions.
% lean: blockScheduleModel
\end{enumerate}
\end{proof}
